% !TEX root = ../main.tex

% ====

\section{Security Proof}

\subsection{Overview}

Zeeperio is a component of a voting system, not a full-fledged voting system. We will keep our proof very narrow to the changes Zeeperio introduces because security arguments for most of the other components already exist. Eperio is a `back-end' tabulator that proves the correctness of the election, while leaking no information about how any voter voted from all its transcripts. The `front-end' of the system consists of the paper ballots, how they encode votes, and how voters interact with their voted ballots, audited ballots, and the transcript of the election. The front-end also needs to be receipt-free. Eperio, and by extension Zeeperio, is compatible with multiple front-ends, but we use the example of Scantegrity II.

Scantegrity II is an end-to-end verifiable voting system with paper ballots that are optically scanned. A ballot contains a unique serial number and it contains a unique (per ballot) confirmation code printed beside each candidate in invisible ink. When the voter marks a candidate on the ballot, the voter is able to also read the code. The ballot serial number and the revealed code constitutes the voter's receipt, which they can copy to a piece of paper and bring home. The scanned ballots are provided to a backend system called `the switchboard.' The switchboard prepares a proof that the tally is correct and provides an interface for voters to see if their receipts are recorded correctly and if their print audited ballots contain consistent confirmation codes with what is recorded.

\subsection{Scantegrity II is Coercion Resistant}

For the purposes of this paper, we will assume that Scantegrity II is a secure voting system. Its security is argued in the Scantegrity papers, while a very thorough proof of its \textit{coercion resistance} property is given by Kusters, Truerung, and Vogt. The definition of coercion resistance is from their earlier paper at CSF 2010. We provide a summary of their main theorem, while omitting many of the details and exact models (see their paper).

Let $P_{\mathsf{Sct}}$ denote the Scantegrity~II protocol, and let
\[
S = P_{\mathsf{Sct}}(k,m,n,\tilde p)
\]
be an election system with the following parameters. $k$ is the number of candidates (excluding abstention), indexed by $i\in\{1,\dots,k\}$. $m$ is the number of dishonest voters whose votes are assumed known to the coercer (static adversary). $n$ is the number of honest voters (excluding one coerced voter). $\tilde p=(p_0,\dots,p_k)$ is the distribution of an honest voter's choice ($p_0$ for abstention, $p_i$ for candidate $i$), with $\sum_{j=0}^k p_j=1$.


Kusters \etal define coercion-resistance roughly as: for any strategy $v$ the coercer forces on the voter, there exists a counter-strategy $\tilde v$ such that (i) the resulting run satisfies the voter's goal $\gamma_i$ with overwhelming probability, while (ii) the coercer cannot distinguish whether the voter obeyed the coercer or used its counter-strategy by more than a security parameter $\delta$. The quantity $\delta=\delta^i_{\min}(n,k,\tilde p)$ is the paper's expression for the best-possible distinguishability advantage available to a coercer in this setting (we refer the interested reader to the paper for its exact details).

\begin{theorem}[Main Result from TK]
Fix $i\in\{1,\dots,k\}$ and let $\delta \;=\; \delta^i_{\min}(n,k,\tilde p).$ Then the Scantegrity~II election system $S=P_{\mathsf{Sct}}(k,m,n,\tilde p)$ is
$\delta$-coercion-resistant with respect to the goal $\gamma_i$.
\end{theorem}


Under this definition, it does not say ``Scantegrity II is coercion resistant'' writ large. Rather it says, ``Scantegrity II is coercion-resistant for specific family of voter goals $\gamma_i$.'' In general, a system may be proven coercion resistant for certain kinds of goals and not for other goals. For the authors' proof of Scantegrity II, they need a specific goal for the proof. The goal they use, $\gamma_i$, can be informally stated as: \textit{if the coercer’s demand is abstention, the voter abstaining is acceptable; otherwise, if the voter casts a vote, it should be for candidate $i$, regardless of which candidate the coercer instructs them to vote for.}


\subsection{Replacing the Switchboard with Eperio} 

As mentioned, the verification back-end in Scantegrity II is referred to as the Switchboard. Eperio was introduced by its authors explicitly as a replacement for the Switchboard. Eperio tables play the same architectural role as the switchboard, while providing proofs that are asserted to be more human-understandable. The paper shows an Eperio proof is simpler to implement (smaller data, less lines of code) than Switchboard proofs and can even be manually checked by a human verifier using a spreadsheet (such as Microsoft Excel) and simple operations.

A blackbox equivalence of Eperio and the Switchboard has not been formally proven. Basin \etal do not seek to prove this, but they do model both Eperio and the Scantegrity switchboard. They show that they operate over the same underlying table structures and for every verification property they prove, they just state the proof would be analogous to Eperio's proof without developing them further. The also show that the same privacy argument applies to both systems. The privacy result they obtain for Eperio, however, comes with two important caveats relative to the Kusters \etal result. First, the notion of privacy is weaker than coercion resistance: it shows that the published verification data does not allow votes to be linked to voters within a table-based model, a property sometimes referred to as receipt freeness, rather than a full game-based coercion-resistance guarantee. Second, their analysis is framed using a simplified Scantegrity-style ballot abstraction rather than the full Scantegrity II ballot format, so the result applies at the level of the switchboard and verification logic, not to the complete human-facing voting protocol.

We state as a conjecture, and without proof, that if the Switchboard is replaced with Eperio, the result from Kuster \etal will continue to hold. We think this is reasonable given the evidence above.

\subsection{Zeeperio is as Coercion Resistant as Eperio}

% !TEX root = ../main.tex
\clearpage

We now turn to Zeeprio and begin with coercion resistance. The key observation we call non-expansivity: Zeeprio does not reveal any information beyond what is already revealed by Eperio. More precisely, any information available to a coercer from a Zeeprio transcript can be derived from an Eperio transcript, possibly with weaker assurance. Consequently, Zeeprio cannot be less coercion resistant than Eperio: any privacy property that holds for Eperio continues to hold when augmented with Zeeprio’s additional verification artifacts. 

In particular, if non-expansivity holds, Zeeprio inherits at least the level of privacy established for Eperio by Basin \etal~\cite{BDGR21}, namely receipt-freeness under their table-based definition. Moreover, Zeeprio may also inherit coercion resistance in the stronger sense of Küsters \etal~\cite{KTV10c}, provided one accepts the conjecture that replacing Scantegrity II’s switchboard with Eperio does not reduce privacy in any meaningful way.

We formalize this in the following lemma:

\begin{lemma}[Non-Expansive Privacy under Transcript Augmentation]
Let $\Pi$ be a protocol and let $\Pi^{+}$ be an augmentation that, on the same underlying execution, produces a coercer view
$T^{+} := (T,Z)$, where $T$ is the coercer view in $\Pi$ and $Z$ is an additional transcript component.
Assume that for every probabilistic polynomial-time distinguisher $\mathcal{D}$ and every security parameter $\lambda$,
there exists a probabilistic polynomial-time algorithm $\mathsf{Sim}$ such that
\[
\Bigl|\Pr[\mathcal{D}(T,Z)=1] - \Pr[\mathcal{D}(T,\mathsf{Sim}(T))=1]\Bigr|
\]
is negligible as a function of $\lambda$.
\end{lemma}

Under this lemma, $\Pi^{+}$ is not less private than $\Pi$ for any privacy notion whose experiment depends on the protocol only through
the coercer view. In particular, for any such privacy experiment and any probabilistic polynomial-time adversary
$\mathcal{A}$ interacting with $\Pi^{+}$, there exists a probabilistic polynomial-time adversary
$\mathcal{B}$ interacting with $\Pi$ such that the difference between the success probabilities of
$\mathcal{A}$ and $\mathcal{B}$ in their respective experiments is negligible in $\lambda$.

 In the special case where $Z$ is efficiently computable from $T$ (and public parameters), the simulator
$\mathsf{Sim}$ may be chosen to simply compute $Z$ from $T$. In other words, $\mathsf{Sim}$ is pure post-processing. This will be the dominant case for Zeeperio.

To reason about the non-expansive privacy of Zeeprio, we separate two concerns: (i) the \emph{cryptographic} claim that the proof a polynomial vanishes is zero-knowledge and proves nothing beyond the statement that the polynomial vanishes; and (i) the \emph{semantic} claim that the constraints implied by Zeeperio are non-expansive over an Eperio transcript.

Starting with the \emph{cryptographic} claim, we directly assume the PIOP used by Zeeprio satisfies the standard zero-knowledge property: its proof transcript reveals no information beyond the fact that the asserted algebraic relations are true. Concretely, as formalized for PLONK-style polynomial IOPs, there exists a simulator that, given only the statement that the constraints are satisfied, produces a transcript that is indistinguishable from an honestly generated one. This simulator construction is explicit in the PIOP literature~\cite{Sef24,Lip25} and does not depend on the particular semantics of the underlying constraints, so we do not reprove anything regarding this.

We do however add two important implementation details that allow the simulation to go through. The first is that the polynomial commitment scheme needs to be hiding. When using KZG commitments~\cite{KZG10}, it is sufficient to use their second construction informally referred to as the Pedersen-variant of KZG. The second is that the PIOP proof will open committed polynomials at random challenge points. If the adversary has a guess as to the specific polynomial under commitment, it can use this random challenge to confirm its guess. To prevent this, polynomials should incorporate extra random points (off-domain) at interpolation time, where the number of points exceed the number of points that will be eventually opened on the given polynomial. This lacks a standard name in the literature and we propose calling it the MBKM heuristic~\cite{MBKM19} after it was popularized in the zk-SNARK system Sonic.

We spend the rest of the proof dealing with the \emph{semantic} claim: the algebraic relations asserted by Zeeprio’s constraint system encode no more information than what is already revealed by Eperio’s verification logic.


\clearpage



\subsection{Zeeperio is as Verifiable as Eperio} 

Basin \etal model Eperio using tables and SP-algebra (selection–projection table algebra) to state verification conditions and prove that these conditions imply integrity (and related) properties through pen-and-paper formal-methods proofs (as opposed to machine-checked verification). In this section, we show the equivalence of Zeeperio's constraints and their model, which we will call the BDGR model or just BDGR. 

The notation we use the main body of our paper and the notation used by BDGR is largely distinct but sometimes a few symbols overlap. Rather than rewriting the BDGR model in our notation, we keep their notation for easier cross-referencing with their paper but we make it easier for the reader to know we are discussing BDGR by placing the main equations in blue boxes, while using pink boxes for Zeeperio.

% Symbols used in the definition of $P_{\mathrm{con}}$.

In BDGR, the election authority reveals a set of Eperio tables but shows only parts of the tables, never the entire table. The entires in the partial table that are unknown are marked with symbol $\bot$. Each artefact $P\in \mathcal{P}_{\mathrm{auth}}$ is in the set of everything published ($\mathcal{P}_{\mathrm{auth}}$). Iterating over each $P$, the model then iterates over every possible $T \in \mathcal{T}_\bot$ from the set of all tables (full or partial). If a table $T$ is consistent with the revealed $P$, which is to say it is the same in possibly permuted order as $P$ with $\bot$ symbols, it is collected. The intersection of all collected tables over each $P$ is denoted $\bigcap_{P\in \mathcal{P}_{\mathrm{auth}}}\{T\in\mathcal{T}_\bot|P\stackrel{\bot}{=}\pi_{\mathrm{sort}(P)}T\}$. 

In their notation, the set of tables that are consistent with the published tally $t$ is denoted $\mathrm{Ind}^{qu}_t(\pi_S\sigma_{M=1})$ which we will not unpack further here, referring the interested reader to their paper. Together, $\mathcal{P}_{\mathrm{con}}$ is a set of full Eperio tables that are consistent with both the tally and the data published in the audit trail. 

\begin{mdframed}[linewidth=0.6pt,backgroundcolor=verylightblue,innerleftmargin=6pt,innerrightmargin=6pt,innertopmargin=4pt,innerbottommargin=4pt]
\[
\footnotesize
\begin{aligned}
\mathcal{P}_{\mathrm{con}}&:=\mathrm{Ind}^{qu}_t(\pi_S\sigma_{M=1}) \cap \bigcap_{P\in \mathcal{P}_{\mathrm{auth}}}\{T\in\mathcal{T}_\bot|P\stackrel{\bot}{=}\pi_{\mathrm{sort}(P)}T\}
\end{aligned}
\]
\end{mdframed}

Everything a voter $v$ sees across an election is specified as their view of the election $R_v$. The collection of all voters' views is defined as Views.

The BDGR model defines an operation $\owedge$ that works like a key-based merge over a set of tables, allowing the input tables to have unknown entries $\bot$. The key Key specifies when two rows should be treated as the same row. The output $\owedge^{Key}(\cdot)$ is the set of possible tables that are consistent merges of the input tables. If the input tables contradict each other, the output set will be empty. Because merges may leave some information underdetermined, $\owedge^{Key}(\cdot)$ will likely contain many tables. The operator $\min$ is then used to extract the least-informative consistent merged tables. Eventually the model will argue that the min , to talk about a single election outcome, BDGR does not require a single merged table; instead it requires that all candidate tables consistent with the evidence yield the same tally.

In an Eperio table, each row has a unique identifier under the key U (and each row corresponds to a specific markable region on a ballot). Using the unique identifier $U$ as a key, the minimum tables of the combined voter views are as follows. 

\begin{mdframed}[linewidth=0.6pt,backgroundcolor=verylightblue,innerleftmargin=6pt,innerrightmargin=6pt,innertopmargin=4pt,innerbottommargin=4pt]
\[
\footnotesize
\begin{aligned}
\mathrm{Views}&:=\{R_v|v\in V\}\\
\bar{V}&:=\min(\owedge^{U}(\mathrm{Views}))\\
\end{aligned}
\]
\end{mdframed}


The BDGR model identifies six predicates $\psi_i$ that should evaluate to true if Eperio is verifiable, sound, and complete, which is the security property $\phi$ they attempt to prove.

\begin{mdframed}[linewidth=0.6pt,backgroundcolor=verylightblue,innerleftmargin=6pt,innerrightmargin=6pt,innertopmargin=4pt,innerbottommargin=4pt]
\begin{definition}[Definition 4]
\label{def:bdgr}
Given a non-empty sort $Key$ of unique keys, $P_{\mathrm{auth}}$, $\mathrm{Views}$, and $P_{\mathrm{con}}$ as in Section~3, let
\[
C = \owedge^{Key}(P_{\mathrm{auth}}\cup \mathrm{Views}) \cap P_{\mathrm{con}}.
\]
We say that an election’s outcome $tal$ corresponds to the voters’ intent, and denote this property by
$\varphi_{Key,tal}(P_{\mathrm{auth}},P_{\mathrm{con}},\mathrm{Views})$, if the following two conditions are satisfied:
\[
(1)\;\; C \neq \emptyset
\qquad\text{and}\qquad
(2)\;\; \forall T \in C,\; \forall T' \in \min(\owedge^{Key}\mathrm{Views}) : tal(T) = tal(T'),
\]
where $tal$ denotes the tally function.
\end{definition}
\end{mdframed}

To relate the BDGR model to Zeeperio's PIOP constraints, we introduce notation that separates two concerns:
(i) the \emph{cryptographic} claim that a verifier acceptance implies the asserted polynomials really vanish (up to negligible error), and
(ii) the \emph{semantic} claim that these vanishing relations faithfully encode Eperio which we will prove via the 6 BDGR predicates.

Let $\mathrm{pub}$ denote the public input to Zeeperio and let $w$ denote the witness. Let $\pi$ denote the PIOP proof transcript (or its non-interactive form), and let $\mathsf{Acc}(\mathrm{pub},\pi)\in\{\mathsf{true},\mathsf{false}\}$ be the verifier’s acceptance predicate.

Zeeperio enforces $\constraints$ algebraic constraints. We write each constraint $c\in\constraints$ as a polynomial identity over a fixed evaluation domain: for every constraint $c$, there is a corresponding \emph{constraint polynomial}: $V_{c}(X;\mathrm{pub},w)$. Intuitively, $V_c$ is constructed so that the constraint $c$ holds exactly when $V_c$ evaluates to $0$ on every point of the domain (or vanishes). If w let $H$ be the evaluation domain, we can define the vanishing predicate as: $\mathsf{Vanish}(f)\;:\Leftrightarrow\;\forall x\in H:\; f(x)=0$. $\mathsf{Acc}(\mathrm{pub},\pi)$ is the event that the verifier accepts the proof, while $\mathsf{Vanish}$ is the underlying mathematical statement that the claimed polynomials vanish on $H$.

We now state the standard soundness guarantee for polynomial IOPs. There exists a soundness error bound $\epsilon$ such that:
\[
\mathsf{Acc}(\mathrm{pub},\pi)\Rightarrow_{\epsilon}\exists w:\bigwedge_j \mathsf{Vanish}\!\left(V_{c_j}(X;\mathrm{pub},w)\right).
\]
In words, except with probability at most $\epsilon$, an accepted proof implies the existence of a witness $w$ for which all constraint polynomials vanish on $H$. Proof of this can be found in the Plonk paper~\cite{GWC19}, which also uses vanishing polynomials to prove knowledge of values that satisfy any arithmetic circuit (and this can be transformed into a general purpose SNARK). While Zeeperio could be defined as a circuit and turned into a SNARK through Plonk, it is more efficient to define custom vanishing polynomials for Eperio directly, creating a special purpose SNARK. 

Next we turn our attention to showing that satisfying the Zeeperio constraints is strong enough to imply the BDGR verification predicates:
\[
\exists w:\bigwedge_j \mathsf{Vanish}\!\left(V_{c_j}(X;\mathrm{pub},w)\right)\Rightarrow \bigwedge_i \psi_i(\mathcal{P}_{\mathrm{auth}},P_{\mathrm{con}},\mathrm{Views}).
\]
This implication is not cryptographic; it is a semantic correspondence argument. Concretely, we interpret each witness polynomial as a column of an underlying table (via evaluation on $H$), and prove that each vanishing constraint enforces the corresponding BDGR table property (e.g., inclusion, uniqueness, tally consistency, and no-stuffing), thereby establishing each $\psi_i$.

\subsection{Aikam see here}

We rely on standard properties of pairing-based polynomial commitments and SNARKs:
\begin{enumerate}
\item[(A1)] \textbf{PCS binding.} The polynomial commitment scheme used by Zeeperio is computationally binding.
For KZG commitments this follows from standard pairing assumptions.
\item[(A2)] \textbf{Knowledge soundness.} If the verifier (smart contract) accepts a Zeeperio proof,
then there exists a witness assignment
satisfying all enforced polynomial identities on~$H$, except with negligible probability.
For Fiat-Shamir transformed PIOP/SNARK systems this is the standard soundness
guarantee in the random-oracle model.
\end{enumerate}

\BDGR{\psi_{0}&:=\owedge^{\{U\}}(\mathcal{P}_{\mathrm{auth}})\cap\bigcap_{P\in\mathcal{P}_{\mathrm{auth}}}\mathrm{shape}_{P}\neq\emptyset}

In Zeeperio, $\mathcal{P}_{\mathrm{auth}}$ is \emph{defined} as $\bot$-filled projections of a single induced election table $E^Z$.
The table $E^Z$ is deterministically computed from the committed witness columns together with the public template.

Recall (from the BDGR discussion of the merge operator $?_{\mathsf{Key}}(\cdot)$) that
$?_{\{U\}}(\cdot)\neq\emptyset$ is an \emph{existential consistency} statement: there exists at least one
completion/merge consistent with all $\bot$-filled inputs, and $?_{\{U\}}(\cdot)=\emptyset$ iff the inputs contradict.


\begin{lemma}[Zeeperio implies $\psi_0$]
\label{lem:psi0-alone-fixed}
Assume the contract accepts $\Pi_{\mathrm{main}}$ and that (A2) holds. 
Then, except with negligible probability, there exists a witness (committed/auxiliary polynomials) consistent with the public commitments such that all enforced Zeeperio polynomial identities hold on~$H$. 
From this witness assignment and the public template, the induced table $E^Z$ over attributes $\{U,M,S\}$is deterministically defined.
Define the authenticated publication tables
\[
P^{UM}\ :=\ \{\langle U:U(i),\,M:M(i),\,S:\bot\rangle \mid 0\le i<P\},
\]
\[
P^{MS}\ :=\ \{\langle U:\bot,\,M:M(i),\,S:S(i)\rangle \mid 0\le i<P\},
\]
and set $P_{\mathrm{auth}}^{Z}:=\{P^{UM},P^{MS}\}$.
Then $\psi_0(P_{\mathrm{auth}}^{Z})$ holds.
\end{lemma}

\begin{proof}
By construction, $P^{UM}$ and $P^{MS}$ are $\bot$-filled projections of the same table $E^Z$ onto the attribute set $\{U,M,S\}$. 
Hence they both have the same $\mathrm{shape}$, so $\mathrm{shape}_{P^{UM}}\cap \mathrm{shape}_{P^{MS}}\neq\emptyset$. 
Therefore $\bigcap_{P\in P_{\mathrm{auth}}^{Z}}\mathrm{shape}_{P}\neq\emptyset$.

Now, consider $\owedge^{\{U\}}(P_{\mathrm{auth}}^{Z})$.
Because both tables originate from the same $E^Z$, the $U$-keyed tuples in $P^{UM}$ and the corresponding $M,S$ information in $P^{MS}$ are jointly consistent.
This is because a valid completion is given explicitly by the table $E^Z$ itself, which simultaneously extends both $\bot$-filled projections.
Therefore, the set of joint completions under key $U$ is nonempty, and thus
\[
\owedge^{\{U\}}(P_{\mathrm{auth}}^{Z})\ \cap\ \bigcap_{P\in P_{\mathrm{auth}}^{Z}}\mathrm{shape}_{P}\ \neq\ \emptyset.
\]

\end{proof}

\BDGR{\psi_{1}\ :=\ \bigwedge_{v\in V}\ \bigwedge_{P\in\mathcal{P}_{\mathrm{auth}}}\ 
\bigl(\pi_{U,M}R_v\ \subseteq^{\#}\ \pi_{U,M}P\bigr).}






\BDGR{\psi_{1}&:=\bigwedge_{v\in V}\bigwedge_{P\in\mathcal{P}_{\mathrm{auth}}}\bigl(\pi_{U,M}R_v\subseteq^{\#}\pi_{U,M}P\bigr)}

\BDGR{\psi_{\mathrm{print\text{-}audit}}&:=\owedge^{\{U\}}\bigl(\{\pi_{U,S}(R)\mid R\in\mathrm{Views}\}\bigr)\cap\owedge^{\{U\}}(\mathcal{P}_{\mathrm{auth}})\cap\bigcap_{P\in\mathcal{P}_{\mathrm{auth}}}\mathrm{shape}_{P}\neq\emptyset}


\BDGR{\psi_{\mathrm{tally}}&:=t\in\bigcap_{P\in\mathcal{P}_{\mathrm{auth}}}\mathrm{Ind}^{qu}_{P}\bigl(\{\pi_{S}\sigma_{M=1}\}\bigr)}


\BDGR{\psi_{\mathrm{unique}}&:=\owedge^{\{U\}}\bigl(\{\pi_{V,U}R\mid R\in\mathrm{Views}\}\bigr)\neq\emptyset}




\BDGR{\psi_{\mathrm{no\text{-}stuffing}}&:=\forall P\in\mathcal{P}_{\mathrm{auth}}:\#\sigma_{M=1}\pi_{M}P=\#\sigma_{M=1}\pi_{M}V}
$\quad\text{where }V=\min\bigl(\owedge^{\{U\}}(Views)\bigr).$




\vanishP{
\Bigl(\bigwedge_j \mathsf{Acc}(\mathrm{pub},\pi)&\Rightarrow \Bigl(\bigwedge_i \psi_i\Bigr)\\
\psi_{0}\wedge\psi_{1}\wedge\psi_{\mathrm{tally}}\wedge\psi_{\mathrm{print\text{-}audit}}\wedge\psi_{\mathrm{unique}}\wedge\psi_{\mathrm{no\text{-}stuffing}}&\Rightarrow\varphi
}


We prove Zeeperio's constraints imply Basin's predicates. Basin \etal prove that if these predicates hold, then the security definition $\phi$ which is provided as Definition~\ref{def:bdgr} above, is satisfied. \qed



\clearpage{Security Proof}

